\subsection{Institutional investor portfolios}

\begin{qfdefinition}[Institutional investors: common characteristics, investment policy, stakeholders, liabilities, horizons and liquidity by type, and the legal, regulatory and tax constraints on each]\label{def:institutional-investor-characteristics}
\qffield{Statement} \emph{Common characteristics} (Interpretation, from \citep[chap.~1, pp.~4--28]{s7-ang}): large pools governed by boards or committees acting for others; long and sometimes perpetual horizons; explicit liabilities or spending rules; access to private markets and to internal or external managers; fiduciary and regulatory duties; and agency problems between the owners, the board, staff and external managers. Most decisions follow a written investment policy that sets objectives, risk tolerance, the strategic allocation, ranges, benchmarks, responsibilities and reporting.

\emph{By type}:
\begin{itemize}
\item \emph{Defined benefit pension plans}: stakeholders are active and retired members, the sponsor, its shareholders and, in the US, the Pension Benefit Guaranty Corporation; liabilities are promised benefits measured as an accumulated or projected obligation; horizons follow the plan's duration; liquidity is needed for benefit payments. US plans are governed by ERISA, which treats investment fiduciaries as acting solely for participants (pp.~16--22). Minimum funding of a US single-employer plan is set by the Internal Revenue Code: when assets fall short of the \emph{funding target}, the present value of benefits accrued at the start of the plan year, the sponsor pays the target normal cost plus installments on the open shortfall bases, each year's new base being the shortfall less the present value of installments still due on earlier ones. Ang describes the seven-year amortisation of the Pension Protection Act; for plan years beginning after 31 December 2021 (earlier at the sponsor's election) earlier bases were reduced to zero and each new base is amortised over 15 plan years \citep[secs.~430(a), (c)(3), (c)(8), (d)(1)]{us-internal-revenue-code-26}; \citep[secs.~II.B, III.A, III.B.1]{irs-notice-2021-48}. Multiemployer plans follow separate rules.
\item \emph{Sovereign wealth funds}: owned by governments for stabilisation, saving for future generations, reserves or development; liabilities are spending rules or contingent withdrawals, so liquidity must match the purpose (pp.~4--12). The voluntary Santiago Principles call for a sound legal basis, disclosed funding and withdrawal policies, clear governance, investment on risk-adjusted financial grounds unless other aims are disclosed, and a risk-management framework \citep[GAPP~1, 4, 6, 18--19, 22]{santiago-principles-2008}.
\item \emph{Endowments}: perpetual, supporting a university's budget; much of the money is donor-restricted; spending is usually about 4--5 percent (pp.~23--27). State enactments of the Uniform Prudent Management of Institutional Funds Act govern donor-restricted endowment funds; whether a spending level raises a presumption, and how value is measured, depends on the state. In New York, appropriating more than 7 percent of a fund's value, averaged over quarterly values for at least the five preceding years, creates a \emph{rebuttable} presumption of imprudence for gift instruments executed after the act took effect, unless the instrument permits it. Spending at or below 7 percent creates no presumption of prudence; prudence is judged against the statutory factors \citep[sec.~553(a), (d)]{ny-npc-553}.
\item \emph{Private foundations}: perpetual or time-limited grant makers that, in the US, must distribute about 5 percent of non-charitable assets a year or pay an excise tax \citep[sec.~4942(e)]{us-internal-revenue-code-26}; \citep[p.~24]{s7-ang}.
\item \emph{Banks}: liabilities are deposits and borrowing; capital, leverage and liquidity rules constrain the securities portfolio, which manages liquidity and interest-rate risk (\hyperref[def:financial-institution-features]{financial institutions}).
\item \emph{Insurers}: liabilities are policy claims of known (life) or uncertain (property-casualty) timing; solvency rules and asset--liability matching dominate (\hyperref[meth:insurer-analysis]{insurers}).
\end{itemize}
\emph{Taxes}: pensions, endowments and foundations are generally tax-exempt (foundations pay an excise tax); sovereign funds are usually exempt at home and may benefit from immunity or treaties abroad; banks and insurers are taxable (Interpretation).

\qffield{Practice} \emph{Comparing four investors} (original):
\begin{center}\small
\begin{tabular}{@{}p{0.2\linewidth}p{0.24\linewidth}p{0.22\linewidth}p{0.24\linewidth}@{}}
\toprule
Investor & Liability or payout & Horizon and liquidity & Main constraint \\
\midrule
Harlow Industrial DB plan & PBO 9bn, duration 12 & 12 years; benefit payments & IRC sec.~430 funding; sponsor risk \\
Vardø Fund (SWF) & Budget draws in bad years & Decades; liquidity in crises & Spending rule; politics \\
Elmsworth endowment & 5\% spending, smoothed & Perpetual; capital calls & New York NPC sec.~553; donor restrictions \\
Calder Foundation & 5\% minimum payout & Perpetual & s.~4942; excise tax \\
\bottomrule
\end{tabular}
\end{center}

\qffield{Conditions} US federal law as in force for plan years beginning in 2026 for single-employer pensions and foundations; New York law for the endowment; sovereign funds follow national law and the voluntary Santiago Principles. Ang's book predates the 2021 funding changes.

\qffield{Pitfalls} Treating a sovereign fund as having no liabilities. Assuming every institution with a long horizon can hold illiquid assets.

\qffield{See also} \hyperref[def:investor-types-needs]{investor types}; \hyperref[meth:liability-relative-allocation]{liability-relative allocation}.

\end{qfdefinition}

\begin{qfmethod}[Evaluating defined benefit plan risk: funded status, the sponsor's financial strength, correlation between the sponsor's business and the plan's assets, plan design and workforce characteristics]\label{meth:db-pension-risk}
\qffield{Statement} A plan's risk is borne by the sponsor and its shareholders, the members and, through guarantees, taxpayers; the plan's objective should be to meet promised benefits across economic environments \citep[chap.~1, pp.~16--22]{s7-ang}:
\begin{enumerate}
\item \emph{Funded status}: assets against liabilities, which have three measures. The accounting \emph{PBO} projects final salaries and discounts at high-quality corporate bond rates; the statutory \emph{funding target} of a US single-employer plan counts benefits accrued to date at stabilised segment rates, which stay within 95--105 percent of their 25-year averages through 2030 \citep[sec.~430(d)(1), (h)(2)(C)(iv)]{us-internal-revenue-code-26}; an \emph{economic} liability discounts the promised benefits at near risk-free rates and is larger than actuarial values discounted at expected asset returns (pp.~16--21). Only the statutory shortfall drives minimum contributions: since 2022, each year's new shortfall base is amortised over 15 years, replacing the Pension Protection Act's seven (sec.~430(c)(8)). The new base is not the whole shortfall but the shortfall less the present value of the installments still due on earlier bases (sec.~430(c)(3)); \citep[sec.~III.B.1]{irs-notice-2021-48}, so the open bases together amortise exactly the current shortfall.
\item \emph{Sponsor strength}: contributions compete with the sponsor's investment and payouts; a weak sponsor can less easily absorb shortfalls, and guarantees such as the PBGC's act as a put option that rewards risk taking (pp.~17--18).
\item \emph{Sponsor--plan correlation}: if plan assets (or equity holdings in the sponsor's industry) fall when the sponsor's business weakens, deficits appear when contributions are hardest to make (Interpretation).
\item \emph{Plan design}: final-salary benefits tied to wages behave more like real, equity-sensitive liabilities; accumulated benefits like bonds; lump-sum options and indexation add risk; closed or frozen plans shorten horizons (pp.~13--14, 21--22).
\item \emph{Workforce}: older workforces and more retirees mean shorter duration, larger cash outflows and less tolerance for risk; a falling ratio of workers to retirees raises the burden (p.~23).
\end{enumerate}
\emph{Measures}: hedge ratio (asset dollar duration over liability dollar duration), surplus volatility and surplus at risk relative to the sponsor's market value and cash flow.

\qffield{Practice} \emph{Harlow Industrial} (original; \texttt{institutional-checks.py}, part~1). Plan assets 8bn (55 percent global equity, 45 percent bonds of duration 6), PBO 9bn of duration 12; Harlow is a cyclical manufacturer with a 6bn market value and EBITDA of 1.2bn; 60 percent of members are retired and benefits are final-salary with partial indexation.
\begin{itemize}
\item Funded ratio 88.9 percent; the interest-rate hedge ratio is only 20 percent.
\item Surplus volatility is about 1.0bn a year (12.5 percent of assets); the 95 percent one-year surplus at risk of 1.64bn is 27 percent of Harlow's market value.
\item \emph{Three deficits} (plan year 2026; stylised assumptions, not an actuarial valuation): with a funding target of 8.4bn for accrued benefits, the statutory shortfall is 0.40bn; the accounting deficit is 1.0bn; repricing the 9.0bn PBO 0.75 points lower at Treasury rates with duration 12 gives an economic deficit of about 1.81bn.
\item \emph{Statutory minimum for 2026} (stated assumptions, not Harlow's valuation: valuation date 1 January 2026; one 5 percent effective rate in place of the segment rates; open bases set in 2022--2025 of 0.20, 0.08, 0.03 and 0.05bn, with a shortfall every year since, so none was zeroed; no waiver, prefunding or carryover balances; target normal cost 110m). The earlier bases need installments of 33.03m a year, worth 303.2m at the valuation date, so the 2026 base is $400-303.2=96.8$m, amortised by 15 installments of 8.88m from the valuation date \citep[sec.~430(a), (c)(1)--(3), (c)(8), (g)(2)]{us-internal-revenue-code-26}. Minimum required contribution: target normal cost 110m plus a shortfall charge of 41.91m, \emph{151.91m} at 1 January 2026 (12.7 percent of EBITDA), of which only 41.91m (3.5 percent) repairs the shortfall; the rest pays for benefits earned in 2026. Every installment still due has a present value of exactly 0.40bn. Payments after the valuation date are adjusted for interest at the plan's effective rate \citep[para.~(b)(4)(i)]{cfr-26-1-430j-1}: 151.91m at 1 January 2026 is equivalent, on that ordinary adjustment alone, to 165.12m on the last due date, 15 September 2027. That figure is an illustration, not an amount that would satisfy the minimum. Because the plan had a shortfall in 2025, Harlow must pay quarterly installments of 25 percent of the lesser of 90 percent of the 2026 minimum and 100 percent of the 2025 minimum, due 15 April, 15 July and 15 October 2026 and 15 January 2027 (paras.~(c)(1), (c)(5)--(6); sec.~430(j)(3)). Any part of a September 2027 payment that settles a late installment is first discounted to the installment's due date at the effective rate plus 5 percentage points and only then to the valuation date at the effective rate (para.~(b)(4)(ii)), so a single September 2027 payment would have to exceed 165.12m by an amount that depends on the 2025 minimum, which is not modelled here. Treating the whole 0.40bn as a new base and adding the earlier installments would charge 69.73m instead of 41.91m, amortising the 303.2m twice. If assets reached the funding target, every base would be zeroed and the minimum would be target normal cost less the excess (sec.~430(a)(2), (c)(5)--(6)).
\item \emph{Policy illustration, not law}: a level end-of-year annuity clearing the 1.0bn accounting deficit at 5 percent costs 96.34m a year over 15 years (8.0 percent of EBITDA) or 172.82m over seven (14.4 percent). The equity-heavy portfolio is likely to lose money in the recessions in which Harlow's earnings also fall.
\item \emph{Evaluation and recommendation}: three measures, three answers. Under the stated assumptions the statutory cash for 2026 is 151.91m, of which the deficit-repair part (41.91m, 3.5 percent of EBITDA) is moderate; target normal cost is a running cost of the benefits still accruing, not a burden of the deficit. With Harlow's actual bases, segment rates and target normal cost the total could differ materially, so without its actuarial valuation the true minimum cannot be determined. The accounting deficit (1.0bn) drives reported pension cost, and the economic deficit (1.81bn) with the 1.0bn surplus volatility is what shareholders bear: a statutory schedule that spreads cash over 15 years lowers the required cash, not the risk. Risk remains high on all five counts. Move to 30 percent return-seeking assets and 70 percent long bonds of duration 15, raising the hedge ratio to 78 percent and cutting surplus volatility to 0.41bn; avoid industrial and cyclical overweights. Fund above the minimum on a policy schedule (for example the 15-year accounting annuity) while earnings are strong. Consider buying annuities for retirees once funding reaches about 100 percent.
\end{itemize}

\qffield{Conditions} One rate factor and one equity factor, a 100 basis point annual rate volatility and illustrative correlations. Funding law is that of US single-employer plans for plan years beginning in 2026; the funding target, the earlier bases, target normal cost, the Treasury spread and the single 5 percent effective rate are stylised (actual installments and present values use the segment rates). Installments are valued at the valuation date, the first day of the plan year; quarterly installment amounts and prefunding elections are not computed.

\qffield{Pitfalls} Judging the plan by asset returns alone. Ignoring the sponsor's ability to pay. Discounting liabilities at the expected return on assets. Treating the whole funding shortfall as a new amortisation base while also paying installments on the earlier bases. Calling the statutory minimum the economic cost.

\qffield{See also} \hyperref[meth:liability-relative-allocation]{liability-relative allocation}; \hyperref[meth:post-employment-benefit-reporting]{pension accounting}.

\end{qfmethod}

\begin{qfmethod}[Evaluating institutional investment policy statements and portfolios: defined benefit plans, sovereign wealth funds, endowments and foundations, and balance-sheet management at banks and insurers]\label{meth:institutional-ips-portfolios}
\qffield{Statement} \emph{Evaluating an IPS}: does it state the purpose and the liabilities or spending rule; a return objective consistent with them; risk tolerance in the institution's own terms (surplus risk, spending volatility, drawdowns, tracking error); liquidity needs including capital calls; legal, regulatory and tax constraints; the strategic allocation with ranges and rebalancing rules; benchmarks; and who decides what? Policies that cannot be measured or that conflict with the liability fail (Interpretation; \hyperref[meth:investment-policy-statement]{IPS components}).

\emph{Evaluating portfolios by type}:
\begin{itemize}
\item \emph{DB plans}: does the hedging portfolio match liability duration and inflation sensitivity; is the return-seeking portfolio sized to the sponsor's capacity (\hyperref[meth:db-pension-risk]{pension risk})?
\item \emph{SWFs}: does the allocation fit the purpose? Norway's fund has a 70/30 equity--bond strategic benchmark, rebalanced when actual equity deviates by more than 2 points at a month-end, a 1.25 percent limit on expected tracking error, and a fiscal rule that spends about the expected 3 percent real return \citep[benchmark index; risk management; about the fund]{nbim-gpfg-2026}.
\item \emph{Endowments}: large alternatives allocations helped elite endowments, through skill in selection and timing rather than illiquidity alone \citep[NBER WP 14341, abstract, pp.~1--2]{lerner-schoar-wang-endowments-2008}; the average endowment did worse than a 60/40 portfolio over 2001--2012 \citep[pp.~26--27]{s7-ang}. Spending rules smooth payouts.
\item \emph{Foundations}: the 5 percent minimum payout plus inflation and costs sets a high required return.
\item \emph{Banks and insurers}: manage the duration gap between assets and liabilities, liquidity under regulatory ratios and capital consumption of risky assets (\hyperref[meth:bank-camels-analysis]{bank analysis}; \hyperref[meth:life-insurer-analysis]{life insurers}).
\end{itemize}

\qffield{Practice} (original; \texttt{institutional-checks.py}, parts~2--3).
\begin{itemize}
\item \emph{Elmsworth endowment's spending rule}: after two 8 percent years and a $-25$ percent year, spending at 5 percent of value falls 29 percent (105.3m to 75.0m); a rule of 70 percent last year's spending grown with 2.5 percent inflation plus 30 percent of 5 percent of value falls 7 percent (106.0m to 98.5m), peaking at 6.6 percent of the post-crash value. Elmsworth is a New York not-for-profit corporation whose endowment is donor-restricted under post-enactment gift instruments that set no spending rate. Its statutory test uses quarterly values averaged over five years \citep[sec.~553(d)]{ny-npc-553}; taking the five years before the example at 2,000m, the smoothed rule never exceeds 5.28 percent of that average and pure 5 percent spending 5.24 percent, each measured against the quarterly values of its own asset path (the two rules spend differently, so their averages differ). Neither raises New York's presumption, but staying below 7 percent does not make spending prudent: the board must document the statutory factors. Economically, the smoothed rule better fits a university budget.
\item \emph{Calder Foundation IPS}: 5 percent payout, 2.5 percent inflation and 0.5 percent costs need 8.2 percent nominal; an IPS expecting 6.5 percent implies a 26 percent loss of real value over 20 years. The IPS should either accept a finite life, raise expected return within its risk tolerance, or seek new gifts; it should not set an unattainable target.
\item \emph{Bank Ardent}: assets of 100 with duration 4 against liabilities of 90 with duration 1.5 leave a duration gap of 2.65; a 1-point rate rise cuts equity by 2.55, a quarter of its 10. The ALM committee should shorten assets or lengthen funding with swaps.
\item \emph{An insurer} with assets of 120 (duration 7) and liabilities of 100 (duration 10) has more liability than asset dollar duration; matching needs asset duration 8.33.
\item \emph{GPFG}: on 30 June 2026 the actual benchmark held 72.93 percent equity, outside the 2-point band around 70, which under the published rule triggers rebalancing at the month-end.
\end{itemize}

\qffield{Conditions} Hypothetical institutions except Norway's published parameters (NBIM web pages, 2026).

\qffield{Pitfalls} An IPS whose return target ignores spending and inflation. Copying the endowment model without the skill or liquidity it needs.

\qffield{See also} \hyperref[def:institutional-investor-characteristics]{institutions}; \hyperref[meth:endowment-liquidity-allocation]{endowment liquidity}.

\end{qfmethod}

\begin{qfmethod}[A sovereign wealth fund risk case: financial, environmental and social risks of the strategy, risks of long-term direct investments, and evaluating and improving enterprise risk management]\label{meth:swf-risk-case}
\qffield{Statement} \emph{Financial risks} of a sovereign fund's strategy: market and factor exposure; correlation with the country's own revenue (oil, commodities, the domestic economy); currency mismatch with future spending; liquidity relative to withdrawal needs, which arrive in crises; leverage and concentration; and political risk of raids on the fund \citep[chap.~1, pp.~7--12]{s7-ang}. \emph{Environmental and social risks}: climate change (physical and transition), human-rights and corruption risks in holdings, and reputational risk that can undermine the fund's legitimacy at home and abroad \citep[chap.~3, pp.~71--73]{s7-ang}; \citep[climate action plan]{nbim-gpfg-2026}.

\emph{Long-term direct investments} (real estate, infrastructure, strategic domestic stakes) add illiquidity, valuation risk, concentration, governance and control risk, and host-country political risk. They are managed by co-investment with experienced partners, diversification across assets and countries, governance rights, independent valuation, exit planning and limits on domestic or policy-driven stakes (Interpretation, with GAPP~15, 19, 21).

\emph{Enterprise risk management} should identify, assess and manage all risks, with reliable information and reporting, controls, incentives and codes of conduct (GAPP~22 \citep[GAPP~22]{santiago-principles-2008}); Norway's manager combines a tracking-error limit with historical and hypothetical stress tests, expected shortfall, concentration and factor analysis \citep[risk management]{nbim-gpfg-2026}.

\qffield{Practice} \emph{The Vardø Fund} (original; \texttt{institutional-checks.py}, part~3), USD~100bn, funded by oil revenue of about 8bn a year, with a rule that the budget may draw up to 6bn in bad years. Allocation: 60 percent global equity, 20 percent bonds, 12 percent direct real estate and infrastructure, 8 percent strategic stakes in domestic companies.
\begin{itemize}
\item \emph{Financial risk}: in a joint oil and equity crash the portfolio falls 21.8bn (22 percent) just as inflows stop and the budget draws 6bn. Bonds worth 20.6bn after the shock cover more than three years of draws, so liquidity is adequate; but the domestic stakes lose 3.2bn, 15 percent of the loss from 8 percent of assets, and are tied to the same economy as oil revenue.
\item \emph{Environmental and social risk}: equity holdings in oil, gas and carbon-intensive industries add to the country's own oil exposure; some infrastructure assets are in countries with weak labour standards.
\item \emph{ERM evaluation}: strengths are a written spending rule and stress tests; weaknesses are that domestic stakes are chosen by the finance ministry outside the risk framework, direct assets are valued by the managing partner, and there is no climate scenario analysis.
\item \emph{Recommendations}: move the domestic stakes to a separate policy vehicle or cap them at 2 percent; require independent quarterly valuation of direct assets; add oil-price and climate-transition scenarios to stress testing and a limit on combined oil exposure (fund holdings plus the country's reserves); publish the governance framework and risk limits (GAPP~16--17); and give the risk function a direct line to the board.
\end{itemize}

\qffield{Conditions} The fund and shocks are hypothetical; Norway's practices are cited only as an example of published arrangements.

\qffield{Pitfalls} Assessing a sovereign fund's risk without the country's other wealth. Treating direct investments' smoothed valuations as low risk.

\qffield{See also} \hyperref[meth:risk-identification-modification]{managing risk exposures}; \hyperref[meth:institutional-esg-integration]{ESG}.

\end{qfmethod}

\begin{qfmethod}[Endowment liquidity: tools for managing portfolio liquidity risk, capturing the illiquidity premium, the denominator effect, and recommending allocation actions for liquidity, risk and return needs]\label{meth:endowment-liquidity-allocation}
\qffield{Statement} \emph{Lessons from 2008.} Harvard's endowment, with 55 percent in hedge funds, private equity and real assets and a third of operating revenue from endowment income, fell 22 percent in four months; secondary buyers offered about half of reported value for its private equity, and the university cut budgets \citep[chap.~13, pp.~410--412]{s7-ang}. Reported illiquid returns overstate returns and understate risk (p.~426).

\emph{Capturing the illiquidity premium}: four ways, through a passive allocation to illiquid classes, choosing less liquid securities within a class, acting as a market maker, or dynamic rebalancing at the portfolio level, which is simplest and has the largest effect (pp.~426, 434). Evidence of a premium \emph{across} classes is weak after biases; within classes it is significant (p.~410).

\emph{Liquidity tools} (Interpretation; \hyperref[meth:alternatives-allocation-approaches]{commitment pacing}): a liquidity budget that classifies assets by days to cash; a reserve of liquid assets covering spending and stressed capital calls; committed credit lines; a commitment-pacing model for unfunded commitments; secondary sales as a last resort; rebalancing ranges wide enough to absorb the \emph{denominator effect} (private assets' share rising when public assets fall and appraisals lag); and derivative overlays to rebalance exposure without selling illiquid assets (\hyperref[ex:endowment-derivative-overlay-case]{overlays}).

\qffield{Practice} \emph{Elmsworth University endowment} (original; \texttt{institutional-checks.py}, part~4). USD~2bn: cash 100m, bonds 200m, public equity 600m, hedge funds 400m (quarterly liquidity), private assets 700m with 500m unfunded; spending 5 percent (100m).
\begin{itemize}
\item \emph{Stress}: public equity $-35$ percent, hedge funds $-10$ percent, private NAV $-10$ percent because appraisals lag: total 1,688m.
\item \emph{Needs}: calls of 125m (a quarter of unfunded) and spending of 100m, less distributions of 13m; \emph{sources}: cash 100m, bonds 208m and a quarter of hedge funds (90m) = 398m, 1.87 times net needs.
\item \emph{Denominator effect}: private assets rise from 35.0 to 37.3 percent of the total while their value falls; with unfunded commitments, exposure is 67 percent. A 35 percent cap would stop new commitments for a year or more.
\item \emph{Recommendation}: adopt a liquidity floor of two years' spending plus a stressed year of calls (325m); Elmsworth meets it before the shock, but after paying 225m and receiving 12.6m of distributions it has 185.6m left, so it should rebuild the floor before committing more; slow new commitments to keep unfunded below 20 percent of assets; arrange a committed 150m credit line; rebalance with futures rather than forced sales; and require private investments to earn at least 2 percentage points a year above public equity (a 1.22 times multiple over ten years) to justify the lock-up.
\end{itemize}

\qffield{Conditions} Stress parameters are illustrative; Harvard's figures are as reported by Ang.

\qffield{Pitfalls} Counting hedge funds with gates as liquid. Ignoring unfunded commitments. Selling public equity at the bottom to meet calls.

\qffield{See also} \hyperref[meth:standards-in-manager-selection]{manager selection}; \hyperref[meth:institutional-ips-portfolios]{spending rules}.

\end{qfmethod}

\begin{qfmethod}[ESG considerations in managing long-term institutional portfolios: motives, approaches from exclusion to stewardship, costs of constraints, climate risk and measurement]\label{meth:institutional-esg-integration}
\qffield{Statement} \emph{Why long-term institutions consider ESG}: as universal owners of diversified portfolios they bear system-wide risks such as climate change; their mandates may require it (Norway's mandate expects portfolio companies to align with net zero by 2050, \citep[climate action plan]{nbim-gpfg-2026}); and beneficiaries' preferences confer legitimacy \citep[chap.~3, pp.~71--73, 112]{s7-ang}.

\emph{Approaches}, as defined by CFA Institute, GSIA and PRI: screening (negative, norms-based, best-in-class), ESG integration into analysis, thematic investing, stewardship (engagement and voting) and impact investing \citep[Tables~1--2, p.~7]{sustainable-investment-review-2022}. Norway combines engagement with exclusions recommended by an independent Council on Ethics, as with Wal-Mart in 2006 (Ang, pp.~71--73).

\emph{Costs and evidence}: exclusions shrink diversification, but excluding a few sectors from a diversified portfolio costs little in mean--variance terms; screening on characteristics is an active strategy judged like any other, and the evidence on returns is mixed (pp.~110--112). Fiduciary duty requires that non-financial aims be disclosed (GAPP~19.1 \citep[GAPP~19.1]{santiago-principles-2008}) and that the cost be measured against the full opportunity set (Ang, p.~112).

\emph{Measurement}: tracking error from exclusions, weighted average carbon intensity, scenario analysis of transition and physical risk, and engagement outcomes.

\qffield{Practice} (original; \texttt{institutional-checks.py}, part~5).
\begin{itemize}
\item \emph{Exclusion cost}: removing a sector with 1.5 percent benchmark weight whose return differs from the rest with 25 percent volatility adds about 0.375 percent of tracking error; four sectors totalling 5.5 percent add 0.72 percent, within a 1.25 percent budget.
\item \emph{Carbon}: holdings of 40/35/25 percent with intensities of 20, 150 and 600 tonnes per USD~million of revenue give a weighted intensity of 210.5; moving 10 points from the most to the least intensive cuts it to 152.5.
\item \emph{Recommendation for the Elmsworth endowment}: adopt a policy based on engagement and voting through its managers, a short norms-based exclusion list approved by the board, climate scenarios in the annual stress test, and reporting of tracking error from exclusions; avoid broad exclusions that would push tracking error beyond its budget without a stated mandate.
\end{itemize}

\qffield{Conditions} The carbon and tracking-error figures are illustrative; ESG ratings differ across providers.

\qffield{Pitfalls} Treating ESG integration as a guarantee of higher returns. Excluding without measuring the cost.

\qffield{See also} \hyperref[meth:esg-exposure-evaluation]{ESG exposures}; \hyperref[meth:swf-risk-case]{SWF case}.

\end{qfmethod}
